\section{Zeckendorf 规范形与稳定化核 \texorpdfstring{$\mathrm{Fold}_m$}{Fold\_m}}
\label{sec:zeckendorf}

\subsection{形式化 Zeckendorf 定理}

形式化证明库给出 Zeckendorf 定理的机械检验证明：每个正整数可唯一表示为若干互不相邻的 Fibonacci 数之和 \cite{AFPZeckendorf}。
该唯一性用于支撑本文的“规范形唯一性”。

\subsection{合法语言与计数}

定义
\[
\Omega_m=\{0,1\}^{m},\qquad
X_m^{\mathrm Z}=\{x\in\{0,1\}^{m}:\text{不含子词 }11\}.
\]
组合证明给出无相邻 $1$ 的 Fibonacci 指示串计数满足 Fibonacci 递推，并可导出 $\card{X_m^{\mathrm Z}}=F_{m+2}$ \cite{PlanetMathZeckendorfCombinatorial}。
因此 $\card{\Omega_6}=64$、$\card{X_6^{\mathrm Z}}=F_8=21$。

\begin{remark}
上述 $2^6\to F_8$ 的计数仅作为主文固定 $m=6$ 的最小可全枚举验证平台：它允许对 $\mathrm{Fold}_6$ 的输出类型与纤维结构进行穷举核查，而不引申为额外语义断言。
\end{remark}

\subsection{稳定化核}

定义满射
\[
\mathrm{Fold}_m:\Omega_m\twoheadrightarrow X_m^{\mathrm Z},
\]
将任意窗口词归约为唯一规范形。附录 \ref{app:zeckendorf-fold} 给出一类局域重写构造与终止/唯一性证明纲要，逻辑依据为 Zeckendorf 唯一表示 \cite{AFPZeckendorf}。

